\documentclass[10pt,a4paper]{article}

% -------------------------------------------------
% Encodage et langue
% -------------------------------------------------
\usepackage[utf8]{inputenc}
\usepackage[T1]{fontenc}
\usepackage[french]{babel}
\usepackage{lmodern}

% -------------------------------------------------
% Mise en page
% -------------------------------------------------
\usepackage{geometry}
\geometry{margin=2.5cm}

\usepackage{setspace}
\setstretch{1.2}

\usepackage{microtype}

% -------------------------------------------------
% Mathématiques
% -------------------------------------------------
\usepackage{amsmath,amssymb,amsthm}
\usepackage{mathtools}

% -------------------------------------------------
% Figures et tableaux
% -------------------------------------------------
\usepackage{graphicx}
\usepackage{float}
\usepackage{booktabs}

% -------------------------------------------------
% Cryptographie
% -------------------------------------------------
\usepackage{cryptocode}

% -------------------------------------------------
% Hyperliens et références
% -------------------------------------------------
\usepackage[hidelinks]{hyperref}
\usepackage[capitalize]{cleveref}
\pcfixcleveref

% -------------------------------------------------
% Listes
% -------------------------------------------------
\usepackage{enumitem}

% -------------------------------------------------
% Environnements
% -------------------------------------------------
\theoremstyle{definition}
\newtheorem{definition}{Definition}[section]
\newtheorem{example}[definition]{Example}

\theoremstyle{plain}
\newtheorem{theorem}[definition]{Theorem}
\newtheorem{lemma}[definition]{Lemma}
\newtheorem{proposition}[definition]{Proposition}
\newtheorem{corollary}[definition]{Corollary}

\theoremstyle{remark}
\newtheorem{remark}[definition]{Remark}

% -------------------------------------------------
% Commandes utiles
% -------------------------------------------------
\newcommand{\sample}{\xleftarrow{\$}}

\newcommand{\Z}{\mathbb{Z}}
\newcommand{\G}{\mathbb{G}}
\newcommand{\GT}{\mathbb{G}_T}
\newcommand{\negl}{\mathrm{negl}}

\usepackage{tikz}
\usetikzlibrary{arrows.meta}
\newcommand{\transition}[1]{%
\medskip
\noindent
\begin{minipage}[c]{0.05\linewidth}
\centering
\begin{tikzpicture}
\draw[-{Latex[length=3mm,width=2mm]}, line width=1pt] (0,0.5)-- (0,-0.5);
\end{tikzpicture}
\end{minipage}%
\hfill
\begin{minipage}[c]{0.92\linewidth}
#1
\end{minipage}
\medskip
}

% -------------------- Title Page --------------------
\title{Titre du rapport de stage}
\author{Nom Prénom}
\date{Date}

% -------------------- Document --------------------
\begin{document}

% ----------- Page de garde -----------
\begin{titlepage}
    \centering

    \vspace*{2cm}

    {\LARGE \textbf{Blind Signatures from Oblivious Transfer:\\ A Detailed Analysis}}\\[1.5cm]

    {\large Julian Mouthon}\\[0.5cm]

    {\large LIP6 --- Sorbonne Université}\\[0.5cm]

    {\large Supervisé par Ky Nguyen}\\[1cm]

    {\large\today}



\end{titlepage}





\newpage

\tableofcontents
\newpage




\section{Introduction}
Les signatures aveugles (Blind Signatures) introduites par Chaum en 1982, permettent à un utilisateur
d'obtenir une signature sur un message déterminé sans que le signataire n'apprenne
son contenu. Ces mécanismes peuvent être utilisés en pratique dans le cadre 
des votes électroniques, des transactions bancaires ou des protocoles d'identification.

Ce rapport présente un nouveau schéma de signatures aveugles utilisant le transfert inconscient
(Oblivious Transfer, OT) et les signatures de Waters.
%============================================================

\section{Préliminaires}

Dans cette partie, nous introduisons les notations et concepts de base nécessaires à la compréhension du protocole présenté dans ce rapport.
Le protocole repose principalement sur les signatures de Waters, un schéma de signature basé sur les groupes bilinéaires,
et le transfert inconscient (OT), plus précisément le OT dual-mode, qui permet de réaliser un transfert de messages de manière sécurisée.

\subsection{Notations}

\begin{itemize}
    \item $\lambda$ désigne le paramètre de sécurité.
    
    \item Un algorithme probabiliste en temps polynomial (PPT) est un algorithme
    dont le temps d'exécution est polynomial en $\lambda$.
    
    \item Pour un ensemble fini $S$, la notation
    $
    x \sample S
    $
    désigne un tirage uniforme de $x$ dans $S$.
    
    \item Pour $\ell\in\mathbb{N}$, on note
    $
    [\ell]=\{1,\ldots,\ell\}.
    $
    
    \item Une fonction $f(\lambda)$ est dite négligeable, notée
    $f(\lambda)=\negl(\lambda)$, si elle décroît plus rapidement que
    l'inverse de tout polynôme en $\lambda$.
    
    \item $\bot$ désigne l'élément d'échec retourné par un algorithme.
\end{itemize}

\subsection{Groupes bilinéaires}

La construction repose sur un système de groupes bilinéaires.
Soient $\mathbb{G}_1$, $\mathbb{G}_2$ et $\mathbb{G}_T$ trois groupes
cycliques d'ordre premier $p$. Un couplage bilinéaire est une application :

\[
e:\mathbb{G}_1\times\mathbb{G}_2
\rightarrow
\mathbb{G}_T 
\]
satisfaisant les propriétés suivantes:

\begin{itemize}
    \item \textbf{Bilinéarité :}
    pour tout $g_1\in\mathbb{G}_1$, $g_2\in\mathbb{G}_2$ et
    $a,b\in\mathbb{Z}_p$,
    \[
    e(g_1^a,g_2^b)=e(g_1,g_2)^{ab}.
    \]
    \item \textbf{Non-dégénérescence :}
    si $g_1$ et $g_2$ sont des générateurs,
    alors
    \[
    e(g_1,g_2)\neq 1_{\mathbb{G}_T}.
    \]
    \item \textbf{Efficacité :}
    le calcul du couplage est réalisable en temps polynomial.
\end{itemize}

\subsection{Signature de Waters}

La construction utilise la fonction de hachage de Waters permettant
d'encoder un message binaire dans un élément du groupe $\mathbb{G}_1$.
Soit un message
$
M=(M_1,\ldots,M_\ell)\in\{0,1\}^{\ell}.
$
Les paramètres publics de la fonction sont constitués de
$
u_0,u_1,\ldots,u_\ell\in\mathbb{G}_1 .
$
La valeur associée au message est définie par :
\[
F(M)=u_0\prod_{i=1}^{\ell}u_i^{M_i}.
\]
Ainsi, chaque bit $M_i$ influence la valeur de hachage par
l'intermédiaire du facteur $u_i^{M_i}$.


Le schéma de signature de Waters utilise les paramètres publics
précédents ainsi qu'une clé secrète $a\in\mathbb{Z}_p$.
La clé publique est 
$
\mathsf{bvk}=g_2^a .
$
Pour signer un message $M$, le signataire choisit un aléa
$t\sample\mathbb{Z}_p$ et calcule :
\[
\sigma_1=h_s^a F(M)^t,
\qquad
\sigma_2=g_2^t .
\]
La signature est alors 
\[
\sigma=(F(M),\sigma_1,\sigma_2).
\]
La vérification consiste à vérifier l'équation de couplage 
\[
e(h_s,\mathsf{bvk})\cdot e(F(M),\sigma_2)
\stackrel{?}{=}
e(\sigma_1,g_2).
\]

\subsection{Oblivious Transfer (OT)}

TODO : FAIRE PARTIE SUR OT

\subsection{Signatures aveugles}

Pour rappel, notre protocole vise à permettre à un utilisateur 
d'obtenir une signature sur un message tout en préservant la 
confidentialité de celui-ci vis-à-vis du signataire.

Un schéma de signature aveugle est constitué des algorithmes suivants:

\begin{itemize}
    \item $\mathsf{BSGen}(1^\lambda)$ :
    génère une clé publique $\mathsf{bvk}$ et une clé secrète
    $\mathsf{bsk}$ ;

    \item $\mathsf{BSUser}(\mathsf{bvk},M)$ :
    exécuté par l'utilisateur, produit un premier message
    $\rho_1$ ainsi qu'un état interne $st$ ;

    \item $\mathsf{BSSigner}(\mathsf{bsk},\rho_1)$ :
    exécuté par le signataire, produit une réponse $\rho_2$ ;

    \item $\mathsf{BSDerive}(st,\rho_2)$ :
    permet à l'utilisateur d'obtenir une signature $\sigma$ ;

    \item $\mathsf{BSVerify}(\mathsf{bvk},M,\sigma)$ :
    vérifie la validité de la signature.
\end{itemize}

Une signature est dite valide si
\[
\mathsf{BSVerify}(\mathsf{bvk},M,\sigma)=1.
\]


\section{Le protocole dual-OT}

\subsection{Motivation : du mode OT classique au mode dual-OT}

\subsection{Présentation}

\subsection{Paramètres publics}

\begin{itemize}
    \item un système de groupes bilinéaires
    $(\mathbb{G}_1,\mathbb{G}_2,\mathbb{G}_T,e,p)$,
    où $\mathbb{G}_1$ et $\mathbb{G}_2$ sont des groupes cycliques
    d'ordre premier $p$, munis d'un couplage bilinéaire
    $
    e:\mathbb{G}_1\times\mathbb{G}_2\rightarrow\mathbb{G}_T,
    $
    avec des générateurs
    $g_1\in\mathbb{G}_1$ et $g_2\in\mathbb{G}_2$.

    \item les paramètres de la fonction de hachage de Waters,
    constitués d'un élément aléatoire
    $u_0\in\mathbb{G}_1$
    ainsi que d'un vecteur
    $\mathbf{u}=(u_1,\ldots,u_\ell)\in\mathbb{G}_1^\ell$.
    La fonction de hachage associée est toujours :
    \[
    F(M)=u_0\prod_{i=1}^{\ell}u_i^{M_i}.
    \]

    \item un CRS global pour l'OT dual-mode :
    $
    \mathsf{OT\text{-}CRS}=(g,h,w,\widetilde{w}),
    $
    où les éléments sont choisis uniformément dans le groupe approprié.
    On définit également :
    $
    \eta=\log_g(h),
    $
    lorsque $g\neq 1$.

    \item la clé publique du schéma de signature :
    $
    \mathsf{bvk}=g_2^a,
    $
    associée à la clé secrète :
    $
    \mathsf{bsk}=h_s^a,
    $
    où $h_s\in\mathbb{G}_1$ est choisi aléatoirement.

\end{itemize}

\subsection{Schéma de signature}


\fbox{%
\begin{minipage}{0.95\linewidth}
\centering

\begin{pchstack}[center]

\procedure{$\mathbf{Signer}$}
{
\\
\\
\\
\\
\\
t\sample\mathbb{Z}_p\\
\forall i\in[\ell]:\
(\mathsf{ek}_{i,1},
\widetilde{\mathsf{ek}}_{i,1})
=
\left(
\dfrac{w}{\mathsf{ek}_{i,0}},
\dfrac{\widetilde w}{\widetilde{\mathsf{ek}}_{i,0}}
\right)
\\
\forall i\in[\ell]:
\text{ choisir }\alpha_i
\text{ tel que }
\prod_i\alpha_i=\mathsf{bsk}\,u_0^t
\\
\forall i\in[\ell],\ b\in\{0,1\}:\\
\qquad s_{i,b},s'_{i,b}\sample\mathbb Z_p\\
\qquad C_{i,b}=g^{s_{i,b}}h^{s'_{i,b}}\\
\qquad
D_{i,b}
=
\mathsf{ek}_{i,b}^{s_{i,b}}
\widetilde{\mathsf{ek}}_{i,b}^{s'_{i,b}}
\alpha_i u_i^{tb}
\\
\sigma_2=g_2^t,\qquad
\sigma'_2=g_1^t
\\
\sendmessageright*
{
(C_{i,b},D_{i,b})_{i,b},
\sigma_2,
\sigma'_2
}
}

\pchspace

\procedure{$\mathbf{User}$}
{
\forall i\in[\ell]:
y_i\sample\mathbb Z_p
\\
(\mathsf{ek}_{i,M_i},
\widetilde{\mathsf{ek}}_{i,M_i})
=
(g^{y_i},h^{y_i})
\\
(\mathsf{ek}_{i,1-M_i},
\widetilde{\mathsf{ek}}_{i,1-M_i})
=
\left(
w\cdot\mathsf{ek}_{i,M_i}^{-1},
\widetilde w\cdot
\widetilde{\mathsf{ek}}_{i,M_i}^{-1}
\right)
\\
\sendmessageleft*
{
(\mathsf{ek}_{i,0},
\widetilde{\mathsf{ek}}_{i,0})_{i\in[\ell]}
}
\\
\\
\\
\\
\\
\\
\\
\\
\\
\\
\\
\\
\\
e(\sigma'_2,g_2) \stackrel{?}{=} e(g_1,\sigma_2)
\\
\sigma_{1,i}
=
D_{i,M_i}C_{i,M_i}^{-y_i}
\\
\sigma_1
=
\prod_{i\in[\ell]}
\sigma_{1,i}
\\
t'\sample\mathbb Z_p
\\
\widehat{\sigma}_1
=
\sigma_1F(M)^{t'}
\\
\widehat{\sigma}_2
=
\sigma_2g_2^{t'}
\\
\widehat{\sigma}'_2
=
\sigma'_2g_1^{t'}
\\
\textbf{return }
\sigma=
(\widehat{\sigma}_1,
\widehat{\sigma}_2,
\widehat{\sigma}'_2)
}

\end{pchstack}

\end{minipage}
}

\subsection{Validité de la signature}

\[
\textsf{Le vérifieur accepte} \iff
\begin{cases}
e(\hat{\sigma}_2', g_2) = e(g_1, \hat{\sigma}_2), \\[0.5em]
e(\hat{\sigma}_1, g_2) = e(h_s, \mathsf{bvk}) \cdot e(F(M), \hat{\sigma}_2).
\end{cases}
\]
Première équation (en utilisant la bilinéarité):
\[
\begin{aligned}
e(\widehat{\sigma}_2',g_2)
&=
e(g_1^{t+t'},g_2)\\
&=
e(g_1,g_2)^{t+t'}\\
&=
e(g_1,g_2^{t+t'})\\
&=
e(g_1,\widehat{\sigma}_2).
\end{aligned}
\]
Seconde équation (toujours pas bilinéarité):
\[
\begin{aligned}
e(\widehat{\sigma}_1,g_2)
&=
e(h_s^aF(M)^{t+t'},g_2)\\
&=
e(h_s^a,g_2)\,
e(F(M)^{t+t'},g_2)\\
&=
e(h_s,g_2^a)\,
e(F(M),g_2^{t+t'})\\
&=
e(h_s,\mathsf{bvk})\,
e(F(M),\widehat{\sigma}_2),
\end{aligned}
\]

\subsection{Les branches DH et messy}

Pour chaque position $i$ du message, l'utilisateur construit deux branches.
La première correspond à son vrai bit $M_i$ et est définie par
\[
(\mathsf{ek}_{i,M_i},\widetilde{\mathsf{ek}}_{i,M_i})
=
(g^{y_i},h^{y_i}),
\]
où $y_i$ est choisi aléatoirement.
Comme $h=g^\eta$, on obtient
\[
h^{y_i}=(g^{y_i})^\eta,
\]
c'est-à-dire
\[
\widetilde{\mathsf{ek}}_{i,M_i}
=
\mathsf{ek}_{i,M_i}^{\,\eta}.
\]
Cette branche vérifie donc la relation de Diffie--Hellman et est appelée \emph{branche DH}.
C'est la branche sur laquelle le protocole fonctionne normalement.
La seconde branche est construite à partir de la première :
\[
(\mathsf{ek}_{i,1-M_i},
\widetilde{\mathsf{ek}}_{i,1-M_i})
=
\left(
\frac{w}{\mathsf{ek}_{i,M_i}},
\frac{\widetilde w}{\widetilde{\mathsf{ek}}_{i,M_i}}
\right).
\]
Pour que cette branche soit également une paire DH, il faudrait que
\[
\widetilde{\mathsf{ek}}_{i,1-M_i}
=
\mathsf{ek}_{i,1-M_i}^{\,\eta}.
\]
Or,
\[
\mathsf{ek}_{i,1-M_i}^{\,\eta}
=
\left(\frac{w}{\mathsf{ek}_{i,M_i}}\right)^\eta
=
\frac{w^\eta}{\widetilde{\mathsf{ek}}_{i,M_i}}.
\]
Ainsi, cette égalité serait satisfaite uniquement si
\[
\widetilde w = w^\eta.
\]
Cependant, le CRS est choisi de manière à vérifier
\[
\widetilde w \neq w^\eta.
\]
La seconde branche ne peut donc jamais être une paire DH. C'est une \emph{branche messy}.
En résumé, pour un utilisateur honnête, chaque position contient exactement une branche DH 
(la bonne branche, correspondant au bit du message) et une branche messy (la mauvaise branche). 
Elle ne transporte aucune information utile.


\section{Preuves de sécurité}

\subsection{Lemmes}

\begin{lemma}[Uniformité des branches messy (non-DH)]\label{lem:lossy}
    
Supposons que $g\neq 1$ et que $\eta=\log_g h$. Soit
$(u,v)\in\mathbb{G}_1^2$ une paire \emph{non-DH}, c'est-à-dire telle que
$v\neq u^{\eta}$. Alors la variable aléatoire
\[
\bigl(C,\mathsf{mask}\bigr)
:=
\bigl(g^{s}h^{s'},u^{s}v^{s'}\bigr),
\qquad
(s,s')\sample\mathbb{Z}_p^2,
\]
est uniformément distribuée sur $\mathbb{G}_1^2$. Par conséquent, pour
tout message clair $m\in\mathbb{G}_1$, le chiffré défini par
\[
(C,D)
=
\bigl(g^{s}h^{s'},u^{s}v^{s'}\cdot m\bigr)
\]
est uniformément distribué sur $\mathbb{G}_1^2$, avec une distribution
qui ne dépend pas de $m$.

\begin{proof}
Puisque $g$ est un générateur de $\mathbb{G}_1$, on peut écrire
$u=g^c$ et $v=g^e$ pour des $c,e\in\mathbb{Z}_p$ uniques.
L'hypothèse $v\neq u^\eta$ implique
$e\neq c\eta\pmod p$.
Considérons l'application $\mathbb{Z}_p$-linéaire
\[
\varphi:\mathbb{Z}_p^2\to\mathbb{Z}_p^2,
\qquad
\varphi(s,s')
=
\bigl(s+\eta s',cs+es'\bigr),
\]
dont la matrice est
\[
\begin{pmatrix}
1 & \eta\\
c & e
\end{pmatrix},
\qquad
\det=e-c\eta.
\]
Comme $e-c\eta\neq 0$ dans le corps $\mathbb{Z}_p$,
$\varphi$ est une bijection de $\mathbb{Z}_p^2$.
Puisque $(s,s')$ est uniformément distribué sur $\mathbb{Z}_p^2$,
$\varphi(s,s')$ l'est également.
Enfin, l'application $x\mapsto g^x$ est une bijection de
$\mathbb{Z}_p$ vers $\mathbb{G}_1$ puisque $g$ est un générateur d'un
groupe d'ordre premier. Appliquée à chaque coordonnée, elle transforme
donc $\varphi(s,s')$ en $(C,\mathsf{mask})$, qui est ainsi uniformément
distribué sur $\mathbb{G}_1^2$.
Pour la seconde affirmation, $D=\mathsf{mask}\cdot m$.
Pour tout $m$ fixé, la multiplication par $m$ est une bijection de
$\mathbb{G}_1$. Ainsi, $(C,D)$ est uniformément distribué sur
$\mathbb{G}_1^2$ ; en particulier, sa distribution est la même pour
tout $m$.
\end{proof}

\end{lemma}

\begin{lemma}[Randomisation des signatures de Waters]
\label{lem:randomisation-waters}
Soit
$
\sigma=(\sigma_1,\sigma_2)
=
\left(h_s^a F(M)^t,g_2^t\right)
$
une signature de Waters valide, où $t\sample\mathbb Z_p$.

Pour tout $t'\sample\mathbb Z_p$, la signature
$
\widehat{\sigma}
=
\left(
\sigma_1F(M)^{t'},
\sigma_2g_2^{t'}
\right)
$
est une signature valide sur le même message $M$ et sa distribution
est indépendante de l'aléa initial $t$.

\begin{proof}
On a
\[
\begin{aligned}
\widehat{\sigma}_1
&= \sigma_1F(M)^{t'}\\
&= h_s^aF(M)^tF(M)^{t'}\\
&= h_s^aF(M)^{t+t'},
\end{aligned}
\]
et
\[
\widehat{\sigma}_2
= \sigma_2g_2^{t'}
= g_2^{t+t'}.
\]
Comme $t'\sample\mathbb Z_p$, la valeur $t+t'$ est uniformément
distribuée dans $\mathbb Z_p$. La signature randomisée est donc
distribuée comme une signature de Waters générée avec un nouvel aléa
uniforme, indépendamment de l'aléa initial $t$.
\end{proof}
\end{lemma}

\begin{lemma}[Indistinguabilité sous DDH]
\label{lem:ddh-waters}
Sous l'hypothèse DDH dans $\mathbb G_1$, les distributions obtenues
en utilisant $F(M)$ et $F(\mathbf{0})$ dans la randomisation sont
computationnellement indistinguables.


\begin{proof}
Supposons qu'il existe un adversaire $\mathcal A$ capable de distinguer
ces deux distributions avec un avantage non négligeable. Nous
construisons alors un distingueur DDH $\mathcal B$.

$\mathcal B$ reçoit une instance
\[
(g,g^a,g^b,c),
\]
où $c$ est soit $g^{ab}$, soit un élément uniformément distribué de
$\mathbb G_1$. Il utilise cette instance pour construire les paramètres
de la simulation.

Si $c=g^{ab}$, la distribution obtenue correspond au jeu dans lequel
la valeur $F(M)$ est utilisée. Si $c$ est uniforme, elle correspond au
jeu dans lequel $F(\mathbf{0})$ est utilisée.

Ainsi, un adversaire distinguant ces deux jeux permet à $\mathcal B$
de distinguer une instance DDH réelle d'une instance aléatoire avec
le même avantage, ce qui contredit l'hypothèse DDH.


\end{proof}
\end{lemma}
TODO : FINIR PREUVE INDISTINGUABILITE

\subsection{Unforgeability}

\subsection{Blindness}


\begin{definition}[Blindness, signataire honnête]
Un schéma de signature aveugle est aveugle pour un signataire honnête si
pour tout adversaire PPT $\mathcal A$,
\[
\left|
\Pr\!\left[
\begin{array}{l}
(\mathsf{bvk},\mathsf{bsk})\gets\mathsf{BSGen}(1^\lambda),\
(m_0,m_1)\gets\mathcal A(\mathsf{bvk}),\
b\sample\{0,1\},\\
(\rho_{1,i},st_i)\gets\mathsf{BSUser}(\mathsf{bvk},m_i)
\quad(i\in\{0,1\}),\\
(\rho_{2,b},\rho_{2,1-b})\gets
\mathcal A(\rho_{1,b},\rho_{1,1-b}),\\
\sigma_i\gets\mathsf{BSDerive}(st_i,\rho_{2,i})
\quad(i\in\{0,1\}),\\
\text{si }\exists i,\
\mathsf{BSVerify}(\mathsf{bvk},m_i,\sigma_i)=0:
\ \sigma_0=\sigma_1=\bot,\\
b=\mathcal A(\sigma_0,\sigma_1)
\end{array}
\right]
-\frac12
\right|
=\negl(\lambda).
\]
La seule différence avec la définition du blindness du protocole
précédent est que $\mathsf{bvk}$ n'est plus choisi par $\mathcal A$,
mais généré honnêtement.

\end{definition}

\begin{proof}[Preuve : blindness pour un signataire honnête]


Pour ce protocole, on ne prouve pas le blindness sous clés malveillantes
(comme pour le protocole précédent), mais une version plus faible où
$(\mathsf{bvk},\mathsf{bsk})$ et le CRS
$(g,h,w,\widetilde w)$ sont générés honnêtement par
$\mathsf{BSGen}$.


On veut montrer que, pour deux messages
$m_0,m_1\in\{0,1\}^{\ell}$, la vue du signataire est indistinguable :
$
\operatorname{View}_{\mathcal A}(m_0)
\approx
\operatorname{View}_{\mathcal A}(m_1),
$
où $\mathcal A$ désigne le signataire honnête.
Comme pour la précédente preuve, nous raisonnons par une suite de jeux.

\bigskip

\noindent\textbf{G$_0$.}

Il s'agit de l'exécution réelle du protocole avec le message $m_0$.
Pour chaque position $i$, l'utilisateur construit une branche DH
correspondant au bit $m_{0,i}$ et une branche messy correspondant à
l'autre valeur. Le signataire calcule ensuite les chiffrés des deux
branches et les transmet à l'utilisateur.

\transition{\textit{On remplace les chiffrés des branches messy par des
couples uniformes.}}

\noindent\textbf{G$_1$.}

Dans un CRS honnêtement généré, chaque position possède une branche DH
et une branche messy. Pour la branche messy, le couple
$
(\mathsf{ek}_{i,1-M_i},
\widetilde{\mathsf{ek}}_{i,1-M_i})
$
est une paire non-DH. D'après le lemme~\ref{lem:lossy}, le chiffrement
réalisé sur cette branche est uniformément distribué sur
$\mathbb{G}_1^2$, avec une distribution indépendante du message
chiffré.
On peut donc remplacer, pour chaque position, le chiffré de la branche
messy par un couple uniformément distribué dans $\mathbb{G}_1^2$ sans
modifier la vue du signataire. Ainsi,
\[
\mathrm{G}_0\approx\mathrm{G}_1.
\]

\transition{\textit{On remplace les branches DH par des branches
indépendantes du bit du message.}}

\noindent\textbf{G$_2$.}

Après le jeu précédent, les chiffrés des branches messy sont
uniformes et ne dépendent plus du message. Il reste à traiter les
branches DH (non messy).
Pour une branche DH correspondant au bit $M_i$, on a
\[
(\mathsf{ek}_{i,M_i},\widetilde{\mathsf{ek}}_{i,M_i})
=
(g^{y_i},h^{y_i}),
\qquad
y_i\sample\mathbb{Z}_p.
\]

Le signataire observe les deux branches sans savoir laquelle
correspond au bit $M_i$. Sous l'hypothèse DDH, il ne peut pas distinguer
la paire
$
(g^{y_i},h^{y_i})
$
d'une paire de la forme
$
(g^{y_i},z_i),
$
$
y_i,z_i\sample\mathbb{Z}_p.
$
On peut donc remplacer la branche DH par une paire
indépendante de la branche correspondant au bit du message.
Après cela, les deux branches associées à chaque position
ont la même distribution quel que soit le bit $M_i$. Par conséquent,
la distribution de l'ensemble du premier message envoyé par
l'utilisateur ne dépend plus de $M$.
Ainsi,
\[
\mathrm{G}_1\approx\mathrm{G}_2
\]

\transition{\textit{On applique la réduction DDH de la
section~\ref{lem:ddh-waters} pour remplacer $F(M)$ par
$F(\mathbf{0})$ dans la signature.}}

\noindent\textbf{G$_3$.}

Après le jeu précédent, les branches envoyées par l'utilisateur sont indépendantes du message $M$.
Après la réponse du signataire, le déchiffrement de la branche DH donne la même forme que dans le protocole classique :
\[
\sigma_1=h_s^aF(M)^t,
\qquad
\sigma_2=g_2^t,
\qquad
\sigma'_2=g_1^t.
\]
On applique alors la réduction DDH de la
section~\ref{lem:ddh-waters}. Celle-ci permet de remplacer, sous
l'hypothèse DDH, le terme $F(M)$ par $F(\mathbf{0})$ dans la signature.
Ainsi, la signature simulée ne dépend plus du message $M$ et
\[
\mathrm{G}_2\approx\mathrm{G}_3.
\]


\transition{\textit{On applique la randomisation de Waters (voir section~\ref{lem:randomisation-waters})}}

\noindent\textbf{G$_4$.}

Après avoir reçu la réponse du signer, l'utilisateur effectue la phase de
déchiffrement. Puis, comme dans le protocole précédent, il effectue la
randomisation de Waters en choisissant un aléa
$t'\sample\mathbb{Z}_p$ et en calculant
\[
\widehat{\sigma}_1
=
\sigma_1 F(M)^{t'},
\qquad
\widehat{\sigma}_2
=
\sigma_2 g_2^{t'},
\qquad
\widehat{\sigma}'_2
=
\sigma'_2 g_1^{t'}.
\]
Comme $t'$ est choisi uniformément dans $\mathbb{Z}_p$, la valeur
$t+t'$ est également uniformément distribuée dans $\mathbb{Z}_p$.
Ainsi, la distribution de la signature finale est indépendante du message.
Elle peut donc être remplacée par une signature simulée utilisant un aléa
indépendant du message. Ainsi,
\[
\mathrm{G}_3\approx\mathrm{G}_4.
\]


\transition{\textit{On remplace $m_0$ par $m_1$.}}

\noindent\textbf{G$_5$.}

Dans ce jeu, le message $m_0$ est remplacé par $m_1$. D'après la transition
précédente, la distribution de la signature finale est indépendante du
message. Ainsi,
\[
\mathrm G_4 \approx \mathrm G_5.
\]

\noindent\textbf{Conclusion.}

On a
$
\mathrm G_0
\approx
\mathrm G_1
\approx
\mathrm G_2
\approx
\mathrm G_3
\approx
\mathrm G_4
\approx
\mathrm G_5
$,
par conséquent,
$
\operatorname{View}_{\mathcal A}(m_0)
\approx
\operatorname{View}_{\mathcal A}(m_1),
$
et le protocole satisfait donc le blindness pour un signataire honnête.

\end{proof}


\section{Implémentation et évaluation expérimentale}

\subsection{Objectifs}

\subsection{Environnement expérimental}

\subsection{Implémentation du protocole}

\subsection{Résultats expérimentaux}

\subsection{Analyse des résultats}


\end{document}